\BourbakiExercisesSection

\begin{enumerate}
\def\labelenumi{\arabic{enumi}.}
\tightlist
\item
  Let A be a ring, E a right A-module and F a left A-module. Let f be a function defined on the set S of all finite sequences
\end{enumerate}

\[
((x _ {1}, y _ {1}), (x _ {2}, y _ {2}), \dots , (x _ {n}, y _ {n})) \quad (n \text {   arbitrary })
\]

of elements of \(\mathbf{E} \times \mathbf{F}\) with values in a set \(\mathbf{G}\) and such that:

\[
f ((x _ {1}, y _ {1}), \dots , (x _ {n}, y _ {n})) = f ((x _ {\sigma (1)}, y _ {\sigma (1)}), \dots , (x _ {\sigma (n)}, y _ {\sigma (n)}))\tag{1}
\]

for every permutation \(\sigma \in \mathfrak{S}_n\) ;

\[
\begin{array}{r l} f \left(\left(x _ {1} + x _ {1} ^ {\prime}, y _ {1}\right), \left(x _ {2}, y _ {2}\right), \dots , \left(x _ {n}, y _ {n}\right)\right) & = f \left(\left(x _ {1}, y _ {1}\right), \left(x _ {1} ^ {\prime}, y _ {1}\right), \left(x _ {2}, y _ {2}\right), \dots , \left(x _ {n}, y _ {n}\right)\right); \end{array} \tag {2}
\]

\[
f ((x _ {1}, y _ {1} + y _ {1} ^ {\prime}), (x _ {2}, y _ {2}), \dots , (x _ {n}, y _ {n})) = f ((x _ {1}, y _ {1}), (x _ {1}, y _ {1} ^ {\prime}), (x _ {2}, y _ {2}), \dots , (x _ {n}, y _ {n}));\tag{3}
\]

\[
f ((x _ {1} \lambda , y _ {1}), \dots , (x _ {n}, y _ {n})) = f ((x _ {1}, \lambda y _ {1}), \dots , (x _ {n}, y _ {n}))\tag{4}
\]

for all \(\lambda \in \mathbf{A}\) .

Show that there exists one and only one mapping \(g\) of \(\mathbf{E} \otimes_{\mathbf{A}} \mathbf{F}\) into \(\mathbf{G}\) such that \(f((x_1, y_1), \ldots, (x_n, y_n)) = g\left(\sum_{i=1}^{n} x_i \otimes y_i\right)\) . (Note that if

\[
\sum_ {i} x _ {i} \otimes y _ {i} = \sum_ {j} x _ {j} ^ {\prime} \otimes y _ {j} ^ {\prime},
\]

the difference \(\sum_{i}(x_i, y_i) - \sum_j(x'_j, y'_j)\) in the \(\mathbf{Z}\) -module \(\mathbf{Z}^{(\mathbf{E} \times \mathbf{F})}\) is a linear combination with integer coefficients of elements of the form (2) of no. 1.)

*2. Consider the field \(\mathbf{C}\) of complex numbers as a vector space over the field \(\mathbf{R}\) of real numbers.

\begin{enumerate}
\def\labelenumi{(\alph{enumi})}
\item
  Show that the canonical mapping \(\mathbf{C} \otimes_{\mathbf{R}} \mathbf{C} \to \mathbf{C} \otimes_{\mathbf{C}} \mathbf{C}\) is not injective.
\item
  Show that the two \(\mathbf{C}\) -module structures on \(\mathbf{C} \otimes_{\mathbf{R}} \mathbf{C}\) arising from each of the factors are distinct*.
\end{enumerate}

\begin{enumerate}
\def\labelenumi{\arabic{enumi}.}
\setcounter{enumi}{2}
\tightlist
\item
  \begin{enumerate}
  \def\labelenumii{(\alph{enumii})}
  \tightlist
  \item
    Let \((\mathbf{E}_{\lambda})_{\lambda \in L}\) be a family of right A-modules and F a left A-module. Show that if F is finitely generated, the canonical mapping
  \end{enumerate}
\end{enumerate}

\[
\left(\prod_ {\lambda \in L} E _ {\lambda}\right) \otimes_ {A} F \rightarrow \prod_ {\lambda \in L} (E _ {\lambda} \otimes_ {A} F)
\]

is surjective.

\begin{enumerate}
\def\labelenumi{(\alph{enumi})}
\setcounter{enumi}{1}
\item
  Give an example of a commutative ring A and an ideal m of A such that the canonical mapping \(\mathbf{A}^{\mathbf{N}}\otimes_{\mathbf{A}}(\mathbf{A}/\mathfrak{m})\to(\mathbf{A}/\mathfrak{m})^{\mathbf{N}}\) is not injective (cf.~§1, Exercise 2 (b)).
\item
  Give an example of a commutative ring A such that the canonical mapping \(A^{N} \otimes_{A} A^{N} \to A^{N \times N}\) is not surjective (observe that for given n, the
\end{enumerate}

image of \(\mathbf{N}\) under a mapping of the form \(m \mapsto \sum_{i=1}^{n} u_i(m)v_i(n)\) , where the \(u_i\) and \(v_i\) belong to \(\mathbf{A}^{\mathbf{N}}\) , generate a finitely generated ideal of \(\mathbf{A}\) .

\begin{enumerate}
\def\labelenumi{(\alph{enumi})}
\setcounter{enumi}{3}
\tightlist
\item
  Deduce from (b) and (c) an example where the canonical homomorphism (22) (no. 7) is neither injective nor surjective.
\end{enumerate}

\begin{enumerate}
\def\labelenumi{\arabic{enumi}.}
\setcounter{enumi}{3}
\tightlist
\item
  Let E be a right A-module and F a free left A-module. Show that if x is a free element of E and y an element \(\neq0\) of F, then \(x\otimes y\neq0\) ; if also A is commutative and y is a free element of F, then \(x\otimes y\) is free in the A-module \(E\otimes_{A}F\) .
\end{enumerate}

